\section*{Leitidee}
Differentialgleichungen beschreiben, wie sich Gr��en in Abh�ngigkeit von ihren �nderungsraten entwickeln. Sie bilden das zentrale mathematische Modell f�r dynamische Systeme.

Im linearen Fall greifen sie dabei auf die zuvor entwickelte lineare Struktur von Funktionenr�umen zur�ck: L�sungen k�nnen addiert, skaliert und �berlagert werden. Dieses Superpositionsprinzip erlaubt es, komplexe dynamische Prozesse aus einfacheren Bausteinen zusammenzusetzen.
